\documentclass[10pt,a4paper]{amsart}
\usepackage[margin=19mm]{geometry}
\usepackage{amsmath,amssymb,mathtools}
\usepackage[scr=rsfs,cal=euler]{mathalfa}
\usepackage{xfrac}
\usepackage[unicode,hidelinks,bookmarks=false]{hyperref}
\newcommand{\ie}{i.e.\ }
\newcommand{\cf}{cf.\ }
\newcommand{\resp}{respectively\ }
\newcommand{\Subsection}{\subsection}
\providecommand{\nobreakdash}{\nobreak\hbox{-}}
\setlength{\emergencystretch}{2em}
\multlinegap0pt
\usepackage{amssymb}
\usepackage{float}
\newtheorem{thm}{Theorem}
\newtheorem{prop}{Proposition}
\newtheorem{cor}{Corollary}
\newtheorem{lemma}{Lemma}
\newtheorem{rem}{Remark}
\newtheorem{definition}{Definition}
\newcommand{\m}[1]{\mathbb{#1}}
\newcommand{\mc}[1]{\mathcal{#1}}
\title{Large deviations for the 3D dimer model}
\author{Nishant Chandgotia, Scott Sheffield, Catherine Wolfram}
\date{Source: arXiv:2304.08468v3 (CC BY 4.0)}
\begin{document}
\maketitle

\noindent\emph{Source and licence.} Large deviations for the 3D dimer model. Version arXiv:2304.08468v3 (CC BY 4.0), distributed by arXiv under the Creative Commons Attribution 4.0 International licence, http://creativecommons.org/licenses/by/4.0/, as recorded in the arXiv licence metadata for this exact version and shown on the abstract page https://arxiv.org/abs/2304.08468v3. Full licence notice: \texttt{licenses/arxiv-2304.08468-CC-BY-4.0.txt}.\par\smallskip

\noindent\textit{Editorial scope.} Counted unit: the article's thm 2 (source0.tex lines 1231--1233). The article's complete original proof of it is delivered with it (source0.tex lines 4250--4432, one \texttt{proof} environment of 183 lines, which the article places away from the statement). No mathematics is written, repaired or re-derived: the statement and the proof are the article's own lines, in the article's own notation, and no retained line is substituted or re-flowed.  Retained uncounted as the source-local context the proof needs: lines 1256--1261; lines 1296--1298; lines 2752--2757; lines 2758--2760; lines 3333--3339; lines 4107--4113; lines 4160--4163; lines 4176--4184; lines 4208--4210; lines 4214--4222.  Nothing was omitted from any retained span, so the package carries no omission record.  No retained line contains a citation, so the wrapper carries no bibliography and no bibliography entry.  No retained line contains a figure environment, an image-inclusion command or drawing code, so no image is delivered and the package declares no asset dependency.  Editorial formatting only, in addition: an AMS \texttt{amsart} preamble that restates the article's own declarations and macros used by the retained lines, namely the environments \texttt{thm}, \texttt{prop}, \texttt{cor}, \texttt{lemma}, \texttt{rem}, \texttt{definition} on separate counters, together with the standard packages those lines need.  The retained environments print in this document as Theorem 1, Proposition 1, Corollary 1, Lemma 1, Remark 1, Definition 1 in order of appearance, whereas the article prints the same items with the numbering of its own full document.  The complete native source of the article as distributed in the arXiv e-print for this exact version is bundled unchanged as \texttt{sources/140219/source0.tex}.
\begin{prop}\label{prop: white black imbalance}
	Suppose that $R$ is balanced but not tileable, and that $U\subset R\subset \m Z^3$ is a counterexample to tileability. Then 
	\begin{equation}\label{eq:imbalance_area}
		0 < \text{imbalance}(U) = \text{even}(U) - \text{odd}(U) = \frac{1}{6} \bigg( \text{white}(\partial U) - \text{black}(\partial U) \bigg).
	\end{equation}
\end{prop}

\begin{cor}\label{cor:minimal_tileable}
    A finite balanced region $R\subset \m Z^3$ is tileable if and only if it has no minimal counterexamples. 
\end{cor}

\begin{lemma}\label{lem:constant_order_box_Wass_bound}
Suppose $B\subset R$ is a connected region with piecewise smooth boundary. If $\mu, \nu$ are component measures of tiling or asymptotic flows and $\m W_1^{1,1}(\mu,\nu)<\delta$, then there is a constant $C(B)$ depending only on $B$ such that
\begin{align*}
    \m W_1^{1,1}(\mu\mid_B, \nu\mid_B) < \delta + (C(B)+1) \delta^{1/2}.
\end{align*}
\end{lemma}

\begin{rem}\label{rem:constant_explained}
    The constant $C(B)$ is not hard to understand and control. The term $C(B)\delta^{1/2}$ is bounded by $2$ times the volume of the $\delta^{1/2}$ annulus with inner boundary $\partial B$. 
\end{rem}

\begin{thm}\label{thm: boundary_value_uniformly_continuous}
For any $S\in \m S(R)$, the trace operator 
\begin{align*}
    T(\cdot,S): (AF(R)\cup TF(R),d_W) \to (\mc M^s(R), \m W_{1}^{1,1})
\end{align*}
is uniformly continuous. In particular this holds for $S = \partial R$.
\end{thm}

\begin{thm}[Generalized patching theorem]\label{thm:generalized_patching}
    Fix $c>0$. Let $(R,b)$ be flexible, with $b$ a boundary asymptotic flow which is extendable outside. Let $R_n\subset \frac{1}{n} \m Z^3$ be a sequence of tileable regions approximating $R$ in Hausdorff distance and with boundary values on $\partial R$ converging to $b$ in $\m W_1^{1,1}$. Let $\sigma_n$ be a sequence of tilings of $R_n$. 

    Let $\tau_n\in T_n(R)$ be a sequence of tilings such that $d_W(g,f_{\tau_n}) \to 0$ as $n\to \infty$ for $g\in AF(R,b)$ flexible. Let $\tau_n'$ be $\tau_n$ restricted to a free boundary tiling of $R^c$, and let $R_n^c\subset \frac{1}{n} \m Z^3$ be the cubes covered by $\tau_n'$. 
    
    For $n$ large enough, $R_n\setminus R_n^c$ is tileable. 
\end{thm}

\begin{lemma}\label{lem:general_annulusbound}
    Let $R_n\subset \frac{1}{n} \m Z^3$ be a region of diameter in lattice squares at least $n$ such that $R_n$ are regions approximating a fixed region $R\subset \m R^3$ with $\partial R$ a piecewise smooth surface. Define $A = R_n\setminus R_n^c$ for some $c>0$. Suppose that $S\subset A$ is a monochromatic minimal discrete surface in $\frac{1}{n}\m Z^3$ with connects the inner and outer boundaries of $A$. Then there exist constant $c_1,c_2$ independent of $S$ and $n$ such that $$ c_1 n^2 \geq \text{area}_n(S) \geq c_2 n^2$$
    where $c_2 \sim c^2$.
\end{lemma}

\begin{lemma}[Generalized indenting lemma]\label{lem:generalized_indenting}
Let $A = R_n \setminus R_{n}^c$, and fix $\beta>0$ small. Let $S$ be a minimal monochromatic surface connecting inner and outer boundaries of $A$. There exists $\epsilon<c$ independent of $S$ such that the following hold for $n$ large enough:
    \begin{enumerate}
        \item Let $A_{a,b}$ denote the annulus between layers $\partial R_n^a$ and $\partial R_n^b$. Then $\text{area}_n(S\cap A_{\epsilon,2\epsilon}) < \beta n^2$. 
        \item There exists $a\in (\epsilon, 2\epsilon)$ such that $\text{length}_n(\partial(\partial R_n^a \cap S)) \leq (\beta/\epsilon) n$. 
    \item A ``ribbon surface" $\gamma$ for $S\cap \partial R_n^a$ is a surface connecting $\partial R_n^a$ to $\partial R_n$ with boundary $\partial (S\cap \partial R_n^a)$ on $\partial R_n^a$. The ``{ribbon area}" is the minimal $\text{area}_n$ of a ribbon surface, and is bounded by $2\beta n^2$. 
    \end{enumerate}
    The analogous bounds hold for some $a'\in (c-2\epsilon, c -\epsilon)$.
\end{lemma}

\begin{definition}[Shining light measures]\label{def:shininglight_measures}
Let $g\in AF(R,b)$ and let $\widetilde g$ be an approximation as discussed above. For each $n$, let $\widetilde \tau_n\in T_n(\widetilde R)$ be a tiling produced by the shining light construction with $\widetilde g$, where the periodic tilings in the tubes have maximum period $r$. Fix a large constant $C$ which is a multiple of $r$. We define a \textit{sequence of shining light measures} $\lambda_n$ for $g$ using $\widetilde g$ so that for each $n$, $\lambda_n$ is uniform measure on tilings of the form $(\widetilde \tau_n + x)\mid_R\in T_n(R)$ for $x\in \frac{1}{n}\m Z^3$ with $|x| \leq C$. 
\end{definition}

\begin{lemma}[Shining light flow bound]\label{lemma:shininglight_flowbound}
    Let $D\subset \text{Int}(R)$, and let $g\in AF(R,b)$ be flexible. Let $\widetilde g$ be piecewise constant approximation of $g$ on a tetrahedral mesh $\mc X$ such that $\widetilde g$ inherits the flexible condition, and such that the union of tetrahedra $\mc D\subset \mc X$ covering $D$ is still contained in $\text{Int}(R)$.
    
     Let $S$ be a monochromatic black surface in $\frac{1}{n}\m Z^3$ with boundary $\partial S$ and $\lambda_n$ be a sequence of shining light measures as in Definition \ref{def:shininglight_measures} for $\widetilde g$. Let $\Theta_D$ be the collection of odd cubes adjacent to $S\cap D$. Let $N = \text{area}_n(S\cap D)$ be the number of squares from $\frac{1}{n} \m Z^3$ on $S\cap D$. Then there is a constant $K_D\in (0,1)$ independent of $S$ such that for all $n$ large enough,
     \begin{align*}
         \m E_{\lambda_n}[|\text{flux}(\widetilde v_\tau,S\cap D)|]\geq K_{D} |\Theta_{D}| +O(n^{-1}N) + O(\text{length}_n(\partial(S\cap D)).
     \end{align*}
     The constant $K_D$ depends only on $D$ and $g$. In particular, it is independent of $\widetilde g$, as long as $\widetilde g$ inherits the flexible condition and is constructed on a small enough mesh $\mc X$. 
\end{lemma}

\begin{thm}\label{thm:halls matching 2}
Suppose that $G = (V,E)$ is a finite bipartite graph with bipartition $G = \mc A\cup \mc B$ and $|\mc A| = |\mc B|$. Then $G$ {\em fails to have} a perfect matching if and only if there exists a connected set $U \subset V$ such that $|U\cap \mc A| > |U \cap \mc B|$ but all boundary elements of $U$ (i.e., elements of $U$ that are adjacent to some point in $V \setminus U$) belong to $\mc B$.
\end{thm}

\begin{proof}[Proof of Theorem \ref{thm:generalized_patching}]
    If $A = R_n\setminus R_n^c$ is not tileable, then by Theorem \ref{thm:halls matching 2} there exists a counterexample region $U\subset A$ which has only odd cubes along its interior boundary $S$, but has 
    \begin{align*}
        \text{imbalance}(U) = \text{even}(U) - \text{odd}(U) > 0. 
    \end{align*}
    By Corollary \ref{cor:minimal_tileable}, we can assume that the interior boundary $S\subset \partial U$ is a minimal monochromatic discrete surface in $\frac{1}{n} \m Z^3$. 

    By Lemma \ref{lem:generalized_indenting}, for any $\beta>0$ we can find $\epsilon$ and inner and outer layers $a_+\in (\epsilon, 2\epsilon)$ and $a_-\in (c-2\epsilon, c-\epsilon)$ such that for $a=a_+$ or $a=a_-$,
    \begin{equation}\label{eq:length_generalized}
        \text{length}_n(\partial (S\cap \partial R_n^a)) \leq (\beta/\epsilon)n
    \end{equation}
    and further such that there is a ribbon surface $\gamma$ for $\partial (S\cap \partial R_n^a)$ such that 
    \begin{equation}\label{eq:ribbon}
        \text{area}_n(\gamma) \leq 2\beta n^2. 
    \end{equation}
    We define the middle annulus $A_{\text{mid}} = R_n^{a_+}\setminus R_n^{a_-}$. Let $U_{\text{mid}}=U\cap A_{\text{mid}}$. 
    
    We now choose a compact set $D\subset \text{Int}(R)$. We can assume that $D$ is contained in $A_{\text{mid}}$. We can find a piecewise-constant approximation $\widetilde g$ of $g$ on a tetrahedral mesh $\mc X$ which inherits the flexible condition. We let $\mc D\subset \mc X$ be the collection of tetrahedra $X$ such that $X\cap D\neq \emptyset$. We can assume that the mesh scale of $\mc X$ is small enough so that the union of all $X\in \mc D$ is still contained in $\text{Int}(R)$. 
        
    By analogous pigeonhole principle arguments in the indenting lemmas, we can further assume that $D$ has width in squares from $\frac{1}{n} \m Z^3$ of at least $cn/2$ and has $\text{length}_n(\partial (S\cap D)) = O(n)$.
    
    Let $N = \text{area}_n(S\cap D)$. By Lemma \ref{lemma:shininglight_flowbound}, we can find a sequence of shining light measures $\lambda_n$ satisfying for $n$ large enough, 
    \begin{align*}
        \m E_{\lambda_n}[|\text{flux}(\widetilde v_\tau,S\cap D)|]\geq K_D |\Theta_D|+O(N n^{-1}) + O(\text{length}_n(\partial(S\cap D)). 
    \end{align*}
    By Lemma \ref{lem:general_annulusbound}, $|\Theta_D|\geq N/6 = \text{area}_n(S\cap D)/6\geq ( c_2 /24 ) n^2$. Therefore 
    \begin{align*}
        \m E_{\lambda_n}[|\text{flux}(\widetilde v_\tau,S\cap D)|] \geq \frac{K_D c_2}{24}n^2 + O(n).
    \end{align*}
    In particular, we can sample a tiling $\tau_{\text{SL}}$ from $\lambda_n$ such that 
    \begin{equation}\label{eq:test_tiling_SL}
        |\text{flux}(\widetilde v_{\tau_{\text{SL}}},S\cap D)| \geq \frac{K_D c_2}{24}n^2 + O(n).
    \end{equation}
    Let $U_{\tau_{\text{SL}}}$ be the cubes covered by $\tau_{\text{SL}}$ restricted to $U_{\text{mid}}$. This is a tileable region, so 
    \begin{align*}
        \text{imbalance}(U_{\tau_{\text{SL}}}) = 0. 
    \end{align*}
    Let $U_{\text{mid}}'\subset U$ be $U_{\tau_{\text{SL}}}$ minus even cubes in $A\setminus U$ with a face on $S$ which are connected to $U$ by $\tau_{\text{SL}}$. By Equation \eqref{eq:test_tiling_SL}, 
    \begin{equation}\label{eq:imbalance_mid}
        \text{imbalance}(U_{\text{mid}}') \leq -\frac{K_D c_2}{24}n^2 + O(n).
    \end{equation}
    Let $U_{\text{shell}} = U\setminus U_{\text{mid}}'$. It remains to bound the imbalance in $U_{\text{shell}}$.
    
    Consider the set $\alpha = U\cap \partial A_{\text{mid}}$. 
    For each connected component $\alpha_i$ of $\alpha$, we form a closed surface using a cylinder ribbon surface component $\gamma_i$ of $\gamma$ and corresponding patch $\alpha_i'\subset \partial A$. Let $\alpha'$ be the union of the $\alpha_i'$ components. 

     Let $V_i$ be the region enclosed by $\alpha_i,\gamma_i$, and $\alpha_i'$, and let $V = \cup_i V_i$. 

     The regions $U_{\text{shell}}\setminus V$ are the \textit{leftover regions}. Let $W$ be a connected component of the leftover region. By construction, $\partial W$ intersects at most one of $\partial R_n$ or $\partial R^c_n$. Thus the boundary condition on $\partial W$ comes from only one tiling, either $\tau_n$ or $\sigma_n$. 

     Suppose it comes from $\sigma_n$, i.e.\ that $\partial W\cap \partial R_n\neq \emptyset$ (the version where it comes from $\tau_n$ is identical, we just make a choice for concreteness). Since $\sigma_n$ can be extended to a tiling of all of $R_n$, we can extend $\sigma_n$ to a tiling covering $W$. Let $W_\sigma$ be the region covered by the tiles from $\sigma_n$ which intersect $W$. Clearly $\text{imbalance}(W_\sigma)=0$, and $W_\sigma\cap \partial R_n = W\cap \partial R_n$. If a tile in $W_\sigma$ crosses $\partial W$, then either
     \begin{itemize}
         \item It crosses $\partial W \cap S$, in which case $W_\sigma$ contains an even cube which is not contained in $W\subset U$. 
         \item It crosses $\partial W\cap \gamma$ (recall that $\gamma$ is the ribbon surface). In this case $W_\sigma$ could have an odd cube which is not in $W$. However the number of these added cubes over all components $W$ of $\text{U}_{\text{shell}}\setminus V$ is bounded by $\text{area}_n(\gamma) \leq 4\beta n^2$.
     \end{itemize}
    Therefore 
    \begin{align*}
        \text{imbalance}(U_{\text{shell}}\setminus V) \leq 4 \beta n^2. 
    \end{align*}
    We now show that 
    \begin{align*}
        \text{imbalance}(U_{\text{shell}}\cap V) \leq \text{imbalance}(V) + 4\beta n^2. 
    \end{align*}
    The $4\beta n^2$ term again comes from the ribbon area. We use ideas analogous to those above. Let $Y$ be a component of $V\setminus (U_{\text{shell}}\cap V)$. First note that $Y$ intersects at most one of $\partial R_n$ and $\partial R_n^c$, so its boundary condition comes from only one tiling. 

    Assume $\sigma_n$ is the tiling which defines the boundary condition on $Y$, and extend it to a tiling which covers $Y$. Let $Y_\sigma$ be the region covered by $\sigma_n$ tiles which intersect $Y$. Clearly $\text{imbalance}(Y_\sigma) = 0$. If a tile in $Y_\sigma$ crosses $\partial Y$, then it is in one of two cases: 
    \begin{itemize}
        \item It crosses $\partial Y \cap S$. Since $Y\subset A\setminus U$, in this case $Y_\sigma$ contains an odd cube which is not in $Y$. This makes the imbalance of $Y$ larger. 
        \item It crosses $\partial Y \cap \gamma$ (recall $\gamma$ is the ribbon surface). In this case $Y_\sigma$ could have an even cube which is not in $Y$. However the number of these added cubes over all components $Y$ of $V\setminus (\text{U}_{\text{shell}}\cap V)$ is bounded by $\text{area}_n(\gamma) \leq 4\beta n^2$.
    \end{itemize}
    Therefore in summary,
    \begin{align*}
        \text{imbalance}(U_\text{shell}) \leq 8\beta n^2 + \text{imbalance}(V).
    \end{align*}
    We now relate imbalance to flux. As the proof of Proposition \ref{prop: white black imbalance} (where we relate black and white surface area to imbalance), given a set $V$, we apply the divergence theorem to the reference flow $\widetilde r(e) =1/6$ for all $e$ in $\frac{1}{n}\m Z^3$ oriented even to odd to get that 
    \begin{align*}
        \text{imbalance}(V) = \text{flux}(\widetilde r, \partial V) = \frac{1}{6}\bigg(\text{white}(\partial V)-\text{black}(\partial V)\bigg).
    \end{align*}
    We apply this to the set $V=\cup_i V_i$. By Equation \eqref{eq:ribbon}, the $\gamma_i$ contribute at most $4\beta n^2$. Therefore 
    \begin{equation}\label{eq:imbalance_as_flux_r}
        \text{imbalance}(U_\text{shell})  \leq \sum_i \text{flux}(\widetilde r, \partial V_i) \leq 12\beta n^2 + |\text{flux}(\widetilde r,\alpha) - \text{flux}(\widetilde r,\alpha')| .
    \end{equation}
    For the second inequality, we orient $\alpha,\alpha'$ to always both have inward-pointing normal vector (i.e. inward on $\partial A$ and inward on $\partial A_{\text{mid}}$), meaning one has the opposite normal vector as when we compute flux for $\partial V_i$. This is why we get a minus sign.     
    
    The non-rescaled flow $\widetilde f_\tau$ is the divergence-free version of the pretiling flow $\widetilde v_\tau$; related by the equation $\widetilde f_\tau(e) = \widetilde v_\tau(e) - \widetilde r(e)$ for all edges $e$ oriented even to odd. The rescaled version has $f_\tau = \frac{1}{n^3} \widetilde f_\tau$.
    
    Since the boundary condition on $\alpha$ is given by $\tau_{\text{SL}}$ and the boundary conditions on $\alpha'$ are given by $\tau_n$ on the inner boundary and $\sigma_n$ on the outer boundary, none of the tiles from the corresponding tilings cross $\alpha,\alpha'$ and hence
    \begin{align*}
        \text{flux}(\widetilde v_{\tau_{\text{SL}}},\alpha) = \text{flux}(\widetilde v_\ast,\alpha') = 0,
    \end{align*}
    where $\widetilde v_\ast$ is equal to $\widetilde v_{\tau_n}$ on the inner boundary of $\partial A$ and is equal to $\widetilde v_{\sigma_n}$ on the outer boundary of $\partial A$. Therefore
    \begin{equation}\label{eq:imbalanace_as_flux}
         \text{imbalance}(U_\text{shell})  \leq 12\beta n^2 + |\text{flux}(\widetilde f_{\tau_{\text{SL}}},\alpha) - \text{flux}(\widetilde f_{\ast},\alpha')|,
    \end{equation}
    where $\widetilde f_{\ast} = \widetilde f_{\tau_n}$ on the inner boundary of $\partial A$ and $\widetilde f_{\ast}=\widetilde f_{\sigma_n}$ on the outer boundary of $\partial A$.

    It remains to bound these flux differences, and this is where we use information about the boundary conditions. First note that for any surface $X$ and any tiling $\tau$ of $\frac{1}{n}\m Z^3$, the flux of the non-rescaled $\widetilde f_\tau$ and the rescaled $f_\tau$ are related by:
    \begin{equation}\label{eq:differ_by_n2}
        \text{flux}(\widetilde f_\tau,X) = n^2 \text{flux}(f_\tau,X).  
    \end{equation}
    We have that the \textit{rescaled versions} of the tiling flows $f_{\tau_n}$ and $f_{\tau_{\text{SL}}}$ (rescaled, so without the tildes) converge as $n\to \infty$ to the $g\in AF(R,b)$ given in the theorem statement, that is, 
    \begin{align*}
        \lim_{n\to \infty} d_W( f_{\tau_n},g) = \lim_{n\to \infty}d_W( f_{\tau_{\text{SL}}}, g) =0.
    \end{align*}
    Recall that $T(\cdot, X)$ denotes the trace operator which takes a flow to its restriction to a surface $X$. By Theorem \ref{thm: boundary_value_uniformly_continuous}, for $X$ fixed and any $f_1,f_2\in AF(R)$, given any $\delta>0$ there exists $\delta_1$ such that if $d_W(f_1,f_2)< \delta_1$ then $\m W_1^{1,1}(T(f_1,X),T(f_2,X))<\delta$. Recall also that $T(g,\partial R) = b$, and that we are given that $T(f_{\sigma_n},\partial R)$ converges to $b$ in $\m W_1^{1,1}$. Given these facts, we can choose $n$ large enough to guarantee the following: 
    \begin{itemize}
        \item For the outer boundary $\partial R$, 
        \begin{equation}\label{eq:bound1}
        \m W_1^{1,1}(T( f_{\sigma_n},\partial R), b) < \delta
        \end{equation}
        \begin{equation}\label{eq:bound2}
            \m W_1^{1,1}(T( f_{\tau_{\text{SL}}},\partial R), b) < \delta.
        \end{equation}
        \item For the inner boundary $\partial R^c$,
        \begin{equation}\label{eq:bound3}
            \m W_1^{1,1}(T( f_{\tau_n},\partial R^c),T(g,\partial R^c))<\delta
        \end{equation}
        \begin{equation}\label{eq:bound4}
            \m W_1^{1,1}(T( f_{\tau_{\text{SL}}},\partial R^c),T(g,\partial R^c))<\delta.
        \end{equation}
        \item Finally, let $\partial R_{\text{mid}}=\partial (R^{a_+}\setminus R^{a_-})$ be the piecewise smooth surface approximated by $\partial A_{\text{mid}}= \partial(R_n^{a_+}\setminus R_n^{a_-})$. Then  
        \begin{equation}\label{eq:bound5}
            \m W_1^{1,1}(T( f_{\tau_{\text{SL}}},\partial R_{\text{mid}}),T(g,\partial R_{\text{mid}}))<\delta.
        \end{equation}
    \end{itemize}   

    Recall also that boundary value flows correspond to measures, and note that $\text{flux}(f,X) = T(f,X)(X)$ is the total mass of the measure $T(f,X)$ on $X$. In particular, for a surface $X$ and tiling $\tau$ of $\frac{1}{n}\m Z^3$ and $B\subset X$,
    \begin{align*}
        \text{flux}(f_\tau, B) = T(f_\tau,X)(B).
    \end{align*}
Lemma~\ref{lem:constant_order_box_Wass_bound} applied to measures $\mu,\nu$ supported on a surface $X$ says that if $\m W_1^{1,1}(\mu,\nu)<\delta$, then for any $B\subset X$,  
    \begin{align*}
        \m W_1^{1,1}(\mu\mid_B,\nu\mid_B) \leq \delta + \delta^{1/2} (C(B)+1),
    \end{align*}
    where $\delta^{1/2} C(B)$ is bounded by $2$ times the difference of the area of B and the $\delta^{1/2}$ neighborhood of $B$ within $X$; equivalently, by the area of the annulus of width $\delta^{1/2}$ with inner boundary $\partial B$ (see Remark \ref{rem:constant_explained}). 
    
To use this, we relate $\alpha,\alpha'$ which are contained in the discrete surfaces $\partial A_{\text{mid}}$ and $\partial A$ built out of $\frac{1}{n}\m Z^3$ lattice squares, to $B,B'$ on the piecewise smooth surfaces $\partial R_{\text{mid}}$ and $\partial R\cup\partial R^c$ respectively.

By Equation \eqref{eq:length_generalized}, $\text{length}_n(\partial \alpha) \leq (2\beta/\epsilon)n$. Correspondingly the Euclidean length is bounded as $\text{length}(\partial \alpha)\leq (2\beta/\epsilon)$. Given this, we can cover $\partial \alpha$ with a collection of cubes $\mc C$ with Euclidean side length $3\epsilon$, with $|\mc C|\leq 2\beta/(3\epsilon^2)$. Since the Euclidean width between $\partial A$ and $\partial A_{\text{mid}}$ is less than $2\epsilon$, $\mc C$ also covers $\partial \alpha'\subset \partial A$.

The Hausdorff distances between $\partial A$ and $\partial R\cup \partial R^c$ and between $\partial A_{\text{mid}}$ and $\partial R_{\text{mid}}$ are both bounded by $2/n$. There are corresponding sets $B\subset \partial R_{\text{mid}}$ and $B'\subset \partial R\cup \partial R^c$ which differ from $\alpha, \alpha'$ respectively by Hausdorff distance at most $2/n$. Thus for $n$ large enough, $\mc C$ also covers $B,B'$. Since $\partial R\cup \partial R^c$ and $\partial R_{\text{mid}}$ are piecewise smooth, there is a constant $C'$ such that the area of either surface restricted to one of the cubes in $\mc C$ is at most $C' \epsilon^2$. Since $|\mc C| \leq 2\beta/(3\epsilon^2)$, there is some constant $C$ such that if $\delta^{1/2}\leq\epsilon$, then
\begin{align}
    \label{eq:C(B) bound}\delta^{1/2} C(B) &\leq C \beta\\
    \label{eq:C(B') bound}\delta^{1/2}C(B') &\leq C \beta.
\end{align}
We have that $\text{length}(\partial \alpha)\leq 2\beta/\epsilon$, so the number of $\frac{1}{n} \m Z^3$ lattice points along $\partial \alpha$ is bounded by a constant times $n$ (the constant here depends on $\beta/\epsilon$). Since the Hausdorff distance between $\alpha$ and $B$ is bounded by $2/n$, for any tiling $\tau$ of $\frac{1}{n} \m Z^3$, the flux of $f_\tau$ through a surface is proportional to the number of $\frac{1}{n}\m Z^3$ lattice points on the surface times $\frac{1}{n^2}$. Thus for any tiling $\tau$, since $f_\tau$ is divergence-free,
\begin{align}\label{eq:alphaB bound}
    |\text{flux}(f_\tau,\alpha) - \text{flux}(f_\tau,B)| \leq O(n^{-1}).
\end{align}
By Equation \eqref{eq:C(B') bound}, since $\text{length}(\partial B')\leq C(B')$, also have that $\text{length}(\partial B')\leq C \beta/\epsilon$. Since the Hausdorff distance between $\alpha'$ and $B'$ is bounded by $2/n$, the number of $\frac{1}{n}\m Z^3$ lattice points on a surface between them is also bounded by a constant times $n$ (the constant here depends on $\beta/\epsilon$). Thus we analogously get that for any tiling $\tau$ of $\frac{1}{n} \m Z^3$,
\begin{align}\label{eq:alphaB' bound}
    |\text{flux}(f_\tau,\alpha') - \text{flux}(f_\tau,B')| \leq O(n^{-1}).
\end{align}
Therefore by Lemma \ref{lem:constant_order_box_Wass_bound}, for $\delta$ such that $\delta^{1/2}<\epsilon$, Equations \eqref{eq:bound1},   \eqref{eq:bound3} to relate $f_*$ to $g$ on $B'$, plus Equation  \eqref{eq:C(B') bound} where we determine the constant $C(B')$, and finally Equation \eqref{eq:alphaB' bound} to relate $f_*$ on $\alpha'$ to $f_*$ on $B'$, we get that
 \begin{align}
    |\text{flux}(f_{*},\alpha') - \text{flux}(g,B')| \leq \delta + \delta^{1/2} + C\beta + O(n^{-1}),
\end{align}
where $f_{*}$ is $f_{\sigma_n}$ on the outer boundary of $\partial A$ and $f_{\tau_n}$ on the inner boundary. Similarly, using Equations \eqref{eq:bound2}, \eqref{eq:bound4}, and \eqref{eq:bound5} to relate $f_{\tau_\text{SL}}$ to $g$, plus Equations \eqref{eq:alphaB bound}, \eqref{eq:C(B) bound} for the constant $C(B)$ and to relate $f_{\tau_\text{SL}}$ on $\alpha, B$, the test tiling $\tau_{\text{SL}}$ satisfies analogous bounds on both $\alpha$ and $\alpha'$:
\begin{align}
    |\text{flux}(f_{\tau_{\text{SL}}},\alpha') - \text{flux}(g,B')| &\leq \delta +\delta^{1/2}+ C\beta + O(n^{-1})\\
     |\text{flux}(f_{\tau_{\text{SL}}},\alpha) - \text{flux}(g,B)| &\leq \delta +\delta^{1/2} + C\beta + O(n^{-1}).
\end{align}
Since $g$ is divergence-free and takes values with norm bounded between $-1$ and $1$, and $B,B'$ differ from $\alpha,\alpha'$ by Hausdorff distance bounded by $2/n$, Equation \eqref{eq:ribbon} implies that
    \begin{equation}\label{eq:bound6}
        |\text{flux}(g,B) - \text{flux}(g,B')| \leq 4\beta + O(n^{-1}).
    \end{equation}

Combining Equation \eqref{eq:differ_by_n2} with the above, 
    \begin{align}
        &|\text{flux}(\widetilde f_{\tau_{\text{SL}}},\alpha) - \text{flux}(\widetilde f_{\ast},\alpha')| \\
        &\leq \bigg[|\text{flux}(f_{\tau_{\text{SL}}},\alpha)-\text{flux}(g,B)| +|\text{flux}(g,B)-\text{flux}(g,B')| +|\text{flux}(g,B')-\text{flux}(f_{*},\alpha')|\bigg] n^2 \\
        &\leq (2 \delta + 2\delta^{1/2} + 2 C \beta + O(n^{-1}) + 4\beta )n^2.
    \end{align}

    Combining this with Equation \eqref{eq:imbalance_mid} and Equation \eqref{eq:imbalanace_as_flux}, we get that
    \begin{align*}
        \text{imbalance}(U) &= \text{imbalance}(U_{\text{mid}}') + \text{imbalance}(U_{\text{shell}}) \\
        &\leq -\frac{K_D c_2}{24}n^2 + (16+2C)\beta n^2 + 2\delta n^2 + 2\delta^{1/2} n^2 + O(n).
    \end{align*}
    The factor $K_D c_2/24$ and the constant $C$ are fixed independent of $\delta,\beta$. Implicit here is also the parameter $\epsilon$ that we indent by.
    
    Taking $\epsilon$ small enough, we can make $\beta$ as small as needed. The parameter $\delta>0$ is related to the distance between tiling flows and their limits and is required to satisfy $\delta^{1/2}<\epsilon$, but this can be guaranteed for $n$ large enough. Therefore for $n$ large enough, $\text{imbalance}(U)$ will be non-positive, and hence $U$ is not a counterexample. This completes the proof.
\end{proof}

\end{document}
